\documentclass{article}

\usepackage{iclr2027_conference,times}
\usepackage[T1]{fontenc}
\usepackage{lmodern}
\usepackage{amsmath,amssymb,amsthm,mathtools}
\usepackage{booktabs,microtype,url,xcolor}
\usepackage[colorlinks=true,allcolors=blue]{hyperref}

\newtheorem{theorem}{Theorem}
\newtheorem{lemma}[theorem]{Lemma}
\newtheorem{proposition}[theorem]{Proposition}
\newtheorem{corollary}[theorem]{Corollary}
\theoremstyle{definition}
\newtheorem{definition}[theorem]{Definition}
\theoremstyle{remark}
\newtheorem{remark}[theorem]{Remark}

\newcommand{\E}{\mathbb E}
\newcommand{\Pp}{\mathbb P}
\newcommand{\cH}{\mathcal H}
\newcommand{\cA}{\mathcal A}
\newcommand{\cY}{\mathcal Y}
\newcommand{\cB}{\mathcal B}
\newcommand{\DeltaS}{\Delta}
\newcommand{\1}{\mathbf 1}
\newcommand{\Reg}{\operatorname{Reg}}
\newcommand{\KL}{\operatorname{KL}}
\newcommand{\TV}{\operatorname{TV}}
\newcommand{\Ber}{\operatorname{Ber}}
\newcommand{\risk}{\mathcal R}

\title{One Extra Label Makes Multiclass Aggregation Mixable}
\author{\textbf{Pahan Dewasurendra} \\ Johns Hopkins University}

\begin{document}
\maketitle

\begin{abstract}
Zero-one classification is not mixable. Aggregating \(N\) multiclass experts therefore costs order \(\sqrt{T\log N}\) regret in general. We show that one extra prediction slot changes the game discontinuously. If every expert predicts at most \(r\) labels and a randomized learner predicts exactly \(k>r\), then expected regret is independent of the horizon.

We characterize the phenomenon exactly. For an unrestricted alphabet, the sharp resource-augmented mixability constant is the positive solution \[ \frac{e^\eta-1}{\eta}=\frac{k}{r}. \] For a pool of \(N\) experts, the sharper constant solves \[ \frac N\eta\log\left(1+\frac{e^\eta-1}{N}\right)=\frac{k}{r}, \] and becomes infinite once \(k\ge Nr\). Exponential weights followed by an explicit exact-cardinality rounding rule then gives regret at most \(\log N/\eta\). We prove matching lower bounds near the phase transition and for larger augmentation.

The same geometry has two learning consequences. It gives a high-probability \(1/n\) excess-risk rate for finite-class agnostic list learning, formalizing a recently identified hedging advantage. It also converts list Littlestone covers into \(O(d\log T)\), rather than square-root-in-\(T\), agnostic regret whenever the learner has more slots than the comparator. Thus a single additional answer changes both adversarial aggregation and stochastic model selection from nonmixable to fast-rate regimes.
\end{abstract}

\section{Introduction}

Top-\(k\) prediction is a basic way to express ambiguity. A classifier may return several diagnoses, recommendations, or candidate tokens and succeeds when the observed label appears in its list. The statistical literature asks how to construct accurate prediction sets or consistent top-\(k\) surrogates \citep{yang2020topk,cortes2024cardinality}. The online literature asks which classes are learnable when predictions are lists \citep{moran2023list}. We ask a more elementary aggregation question. If several experts each return a short list, how much does a learner gain from one additional slot?

Without augmentation, the answer is classical. Randomized zero-one prediction is not mixable, so worst-case regret against \(N\) experts grows as \(\Theta(\sqrt{T\log N})\) \citep{cesabianchi2006prediction}. This remains true if every label is copied \(r\) times and both expert and learner return \(r\)-lists. In contrast, we prove that every strict increase in cardinality makes the asymmetric game mixable. The learner then pays a constant aggregation cost independent of \(T\).

The simplest case is already striking. Experts predict one multiclass label and the learner predicts two. Let \(w_i\) be a posterior over experts and let \(s_y\) be the posterior mass predicting label \(y\). The exponential mix loss at outcome \(y\) is \[ b_\eta(s_y) =-\frac1\eta\log\left(e^{-\eta}+(1-e^{-\eta})s_y\right). \] If the learner includes \(y\) with probability \(p_y\), its expected omission loss is \(1-p_y\). Dominating the mix loss requires \[ p_y\ge g_\eta(s_y) :=1-b_\eta(s_y) =\frac1\eta\log\left(1+(e^\eta-1)s_y\right). \] A randomized two-list exists exactly when these lower marginals use at most two units of inclusion probability. Concavity gives \[ \sum_y g_\eta(s_y) \le g_\eta'(0)\sum_y s_y =\frac{e^\eta-1}{\eta}. \] The largest valid learning rate therefore solves \((e^\eta-1)/\eta=2\), giving \(\eta=1.2564\ldots\). Exponential weights has expected regret at most \(0.796\log N\), regardless of the number of rounds or labels.

This calculation extends exactly in three directions. If experts return \(r\)-lists and the learner returns \(k\)-lists, only the ratio \(k/r\) enters the alphabet-uniform constant. On a finite \(m\)-label alphabet, the exact constant is larger and is attained by the uniform posterior over all \(r\)-subsets. For a fixed pool of \(N\) experts, an even sharper constant is attained by \(N\) disjoint lists. It diverges as \(k\) approaches \(Nr\), at which point the learner can return the union of all expert predictions and obtain pointwise loss no larger than every expert.

The proof is geometric but constructive. Inclusion marginals of randomized exact \(k\)-lists form the hypersimplex \[ \left\{p\in[0,1]^m:\sum_yp_y=k\right\}. \] Mixability reduces to one concave budget inequality. After satisfying its coordinatewise lower bounds, a one-dimensional systematic-sampling rule produces an exact \(k\)-set with the desired marginals. The resulting aggregating algorithm runs in time linear in the union of the expert lists, apart from maintaining the expert weights.

The phase transition is quantitatively sharp. Let \(k=r+s\). A replicated binary hypercube proves an \(\Omega((r/s)\log N)\) horizon-uniform lower bound when \(s/r\) is small. The key identity converts any randomized \(k\)-list into an ordinary binary prediction whose loss is discounted by exactly \(1-s/r\), while comparator losses are not discounted. A noisy Assouad construction then matches the expansion \(\eta=2s/r+O((s/r)^2)\). For larger \(k/r\), a grouped multiclass tree gives \(\Omega(\log N/\log(1+k/r))\), matching the mixability constant whenever the expert pool supports the tree. At \(k=r\), the embedding recovers the usual square-root lower bound.

The result also sharpens two recent learning-theoretic stories. First, \citet{pabbaraju2026optimal} observes that a list predictor may have a hedging advantage when competing with the best single hypothesis and asks whether the usual \(1/\varepsilon^2\) agnostic term is necessary. For a finite class, our answer is a \(1/\varepsilon\) rate. The mix loss is an exp-concave surrogate that agrees with zero-one loss on comparator vertices, and our substitution converts its convex predictions to randomized exact lists. Second, the list Sauer lemma of \citet{moran2023list} covers a class of list predictors by finitely many adaptive list experts. Their generic aggregation incurs a square-root remainder. Resource-augmented mixability reduces it to the logarithm of the cover size.

Our notion is asymmetric. Standard mixability uses the same prediction space for experts and learner \citep{vovk1998game,vanerven2012mixability,reid2015generalized}. Here experts are restricted to \(r\)-lists while the learner receives \(k\)-lists. Standard optimality theorems do not apply because they would also permit experts to use the enlarged action space. This distinction is also why the result is not an instance of evaluating experts by different proper scoring rules \citep{chernov2009evaluators}.

\paragraph{Contributions.}
\begin{enumerate}
\item We define resource-augmented mixability and compute its exact finite-alphabet, finite-expert, and alphabet-uniform constants for list omission loss.
\item We give an efficient aggregating algorithm with horizon-independent regret and a self-contained exact-list rounding rule.
\item We prove lower bounds matching the augmentation dependence near \(k=r\) and away from that threshold.
\item We derive finite-class \(1/n\) stochastic excess risk and logarithmic-in-\(T\) agnostic regret from list Littlestone covers.
\end{enumerate}

\section{Protocol and exact constants}
\label{sec:protocol}

Let \(\cY\) be a label set with at least \(k\) elements. Write \[ \cB_r=\{B\subseteq\cY:|B|\le r\}, \qquad \cA_k=\{A\subseteq\cY:|A|=k\}. \] Both sides are evaluated by omission loss \[ \ell(C,y)=\1\{y\notin C\}. \] Experts choose actions from \(\cB_r\). The learner chooses a distribution over \(\cA_k\).

\begin{definition}[Resource-augmented mixability]
\label{def:mix}
The \((r\!\to\!k)\) list game is \(\eta\)-mixable if, for every finite posterior \(w\) over expert lists \(B_i\in\cB_r\), there is a distribution \(Q\) over \(A\in\cA_k\) such that every \(y\in\cY\) satisfies
\begin{equation}
\label{eq:mix-def}
\E_{A\sim Q}\ell(A,y)
\le-\frac1\eta\log\sum_iw_i e^{-\eta\ell(B_i,y)}.
\end{equation}
\end{definition}

The definition restricts the actions being aggregated. This differs from standard mixability, which quantifies over posteriors on the entire learner action space.

For a posterior \(w\), define its label inclusion marginals \[ s_y=\sum_iw_i\1\{y\in B_i\}. \] Then \(0\le s_y\le1\) and \(\sum_ys_y\le r\). Define
\begin{equation}
\label{eq:bg}
b_\eta(s)=-\frac1\eta\log\left(e^{-\eta}+(1-e^{-\eta})s\right),
\qquad
g_\eta(s)=1-b_\eta(s)
=\frac1\eta\log\left(1+(e^\eta-1)s\right).
\end{equation}
The function \(g_\eta\) is increasing and strictly concave, with \(g_\eta(0)=0\) and \[ g_\eta'(0)=\frac{e^\eta-1}{\eta}. \]

\begin{lemma}[Budget criterion]
\label{lem:budget}
A posterior with inclusion marginals \(s\) admits a substitution satisfying \eqref{eq:mix-def} if and only if
\begin{equation}
\label{eq:budget}
\sum_{y\in\cY}g_\eta(s_y)\le k.
\end{equation}
\end{lemma}

The reason is exact. If \(p_y=\Pp_{A\sim Q}(y\in A)\), then \eqref{eq:mix-def} says \(p_y\ge g_\eta(s_y)\). The feasible marginals of randomized exact \(k\)-sets are precisely the hypersimplex \(p\in[0,1]^\cY\), \(\sum_yp_y=k\). Appendix~\ref{app:geometry} proves both directions without a separation theorem.

We can now state the exact constants. Let \(\eta_m(r,k)\) denote the largest valid constant when \(|\cY|=m\), and let \(\eta_N(r,k)\) denote the largest constant when the posterior is supported on at most \(N\) experts over an arbitrary alphabet. Let \(\eta_\infty(r,k)\) be the largest constant uniform over all alphabet sizes and posterior supports.

\begin{theorem}[Exact resource-augmented mixability]
\label{thm:exact}
Fix \(1\le r<k\).
\begin{enumerate}
\item If \(k<m<\infty\), then \(\eta_m(r,k)\) is the unique positive solution of
\begin{equation}
\label{eq:finite-m}
\frac m\eta\log\left(1+(e^\eta-1)\frac rm\right)=k.
\end{equation}
If \(k\ge m\), then \(\eta_m(r,k)=\infty\).
\item If \(N\ge2\) and \(1<k/r<N\), then \(\eta_N(r,k)\) is the unique positive solution of
\begin{equation}
\label{eq:finite-n}
\frac N\eta\log\left(1+\frac{e^\eta-1}{N}\right)=\frac kr.
\end{equation}
If \(k\ge Nr\), then \(\eta_N(r,k)=\infty\).
\item The alphabet-uniform constant is the unique positive solution
\begin{equation}
\label{eq:infinite}
\frac{e^\eta-1}{\eta}=\frac kr.
\end{equation}
Moreover, \(\eta_m(r,k)\downarrow\eta_\infty(r,k)\) as \(m\to\infty\), and \(\eta_N(r,k)\downarrow\eta_\infty(r,k)\) as \(N\to\infty\).
\end{enumerate}
When \(k=r\), all three nontrivial constants equal zero.
\end{theorem}

The extremizers explain the formulas. On \(m\) labels, uniform mixing over all \(r\)-subsets gives \(s_y=r/m\) and uses budget \(m g_\eta(r/m)\). For \(N\) experts, uniform mixing over \(N\) pairwise disjoint \(r\)-lists gives \(Nr\) active labels with marginal \(1/N\). For an unrestricted alphabet, the worst case is the limit of the first construction as \(m\) grows.

\begin{remark}[Asymptotics]
\label{rem:asymptotics}
Writing \(c=k/r\), the uniform constant obeys \[ \eta_\infty=2(c-1)+O((c-1)^2) \quad\text{as }c\downarrow1, \] and \[ \eta_\infty=\log c+\log\log c+o(1) \quad\text{as }c\to\infty. \] For fixed \(N\), \(\eta_N\sim N\log N/(N-c)\) as \(c\uparrow N\). Thus the finite-pool aggregation cost vanishes continuously when the learner list approaches the union budget.
\end{remark}

\section{Aggregating algorithm}
\label{sec:algorithm}

There are \(N\) experts with prior \(\pi\). On round \(t\), expert \(i\) observes the current instance and returns \(B_{t,i}\in\cB_r\). The learner sees all expert lists, samples \(A_t\in\cA_k\), then observes \(y_t\). The expert predictions may be adaptive.

\begin{minipage}{0.94\linewidth}
\textbf{Resource-augmented aggregating algorithm}
\begin{enumerate}
\item Form posterior weights \(w_{t,i}=W_{t,i}/\sum_jW_{t,j}\) and marginals
\[ s_{t,y}=\sum_iw_{t,i}\1\{y\in B_{t,i}\}. \]
\item Set lower inclusion demands \(d_{t,y}=g_\eta(s_{t,y})\). Increase coordinates without exceeding one until \(p_t\ge d_t\) and \(\sum_yp_{t,y}=k\).
\item Sample an exact \(k\)-list with inclusion marginals \(p_t\) using the systematic procedure in Appendix~\ref{app:rounding}.
\item Update
\[ W_{t+1,i}=W_{t,i} \exp\{-\eta\ell(B_{t,i},y_t)\}. \]
\end{enumerate}
\end{minipage}

Take \(\eta=\eta_N(r,k)\), or the alphabet-uniform \(\eta_\infty(r,k)\) if \(N\) is not used in tuning. Theorem~\ref{thm:exact} guarantees that the demands in Step 2 use at most \(k\) units. Only labels in the union of the expert lists have positive demand. Filler labels can be added in arbitrary order.

\begin{theorem}[Horizon-free regret]
\label{thm:regret}
For every outcome and expert-prediction sequence, every expert \(i\), and every horizon \(T\),
\begin{equation}
\label{eq:regret}
\sum_{t=1}^T\E_t\ell(A_t,y_t)
\le
\sum_{t=1}^T\ell(B_{t,i},y_t)
+\frac{\log(1/\pi_i)}{\eta_N(r,k)}.
\end{equation}
For the uniform prior, expected regret is at most \(\log N/\eta_N(r,k)\). If \(k\ge Nr\), the learner can predict the union of the expert lists on every round and has pointwise loss no larger than every expert.
\end{theorem}

The proof is the standard exponential-weights potential after the nonstandard substitution. The conditional learner loss is at most \(b_\eta(s_{t,y_t})\), which is exactly the one-step change in log potential. Appendix~\ref{app:online} gives the three-line telescope.

The implementation does not enumerate \(k\)-sets. It computes at most \(Nr\) positive demands, fills the remaining budget, and rounds in linear time after sorting labels in any fixed order. Randomization is essential. A deterministic list would need to include every label having positive posterior inclusion, whose number can be \(Nr\).

\section{Lower bounds and the phase transition}
\label{sec:lower}

The exact mixability constants identify the strongest exponential-weights substitution. We next show that their dependence on augmentation is minimax sharp in the principal regimes. Define the horizon-uniform worst-case regret \[ \mathfrak R_N(r,k) =\inf_A\sup_{T\ge1}\sup_{\{B_{t,i},y_t\}} \left\{ \sum_{t=1}^T\E\ell(A_t,y_t) -\min_{i\in[N]}\sum_{t=1}^T\ell(B_{t,i},y_t) \right\}. \]

\begin{theorem}[Augmentation lower bounds]
\label{thm:lower}
There is a universal constant \(c_0>0\) with the following properties.
\begin{enumerate}
\item If \(k=r\), then worst-case regret is \(\Omega(\sqrt{T\log N})\) for the usual range \(2\le N\le 2^T\). In particular, \(\mathfrak R_N(r,r)=\infty\).
\item Let \(k=r+s\) with \(0<s/r\le1/16\). For every \(N\ge2\),
\begin{equation}
\label{eq:near-lower}
\mathfrak R_N(r,r+s)
\ge c_0\frac r s\lfloor\log_2N\rfloor.
\end{equation}
\item Let \(c=k/r>1\), \(q=\lceil2c\rceil\), and \(N\ge q^2\). Then
\begin{equation}
\label{eq:large-lower}
\mathfrak R_N(r,k)
\ge c_0\frac{\log N}{\log(c+1)}.
\end{equation}
\end{enumerate}
\end{theorem}

Since \(\eta_\infty(1+s/r)=\Theta(s/r)\), \eqref{eq:near-lower} matches the alphabet-uniform upper bound \(\log N/\eta_\infty\). Since \(\eta_\infty(c)=\Theta(\log(c+1))\) away from one, \eqref{eq:large-lower} does the same when the pool supports two tree levels. The finite-\(N\) upper bound is sharper as \(c\) approaches \(N\), and zero regret is possible at \(c\ge N\). We do not claim that the grouped-tree lower bound resolves that vanishing finite-pool regime.

\paragraph{Near-threshold idea.}
Use labels \([r]\times\{0,1\}\) and embed a binary expert \(h\) as \[ B_h(x)=\{(j,h(x)):j\in[r]\}. \] Let \(p_{j,b}\) be the learner's inclusion marginals, let \(\bar p_b=r^{-1}\sum_jp_{j,b}\), and put \(a=s/r\). Because \(\bar p_0+\bar p_1=1+a\) and each average is at most one, \[ u_b=\frac{\bar p_b-a}{1-a} \] is a binary distribution. If \(J\) is uniform on \([r]\) and the outcome is \((J,Y)\), then
\begin{equation}
\label{eq:discount}
\E_J[1-p_{J,Y}]
=1-\bar p_Y
=(1-a)(1-u_Y).
\end{equation}
Thus the list learner induces a binary learner whose loss is discounted by \(1-a\), while comparator losses are unchanged.

Take a \(d=\lfloor\log_2N\rfloor\) dimensional binary hypercube with a uniform latent labeling. Query a uniform coordinate and reveal its latent bit through noise with bias \(\gamma=4a\). For \(T=\Theta(d/a^2)\) rounds, a pairwise KL calculation keeps the probability of predicting the latent bit incorrectly bounded below by \(3/8\). By \eqref{eq:discount}, expected per-round regret against the latent expert is at least \[ (1-a)\gamma\frac38-a\frac{1-\gamma}{2} =a+\frac{a^2}{2}. \] The product is \(\Omega(d/a)\). Appendix~\ref{app:lower-proofs} proves the constants and handles arbitrary \(N\).

\paragraph{Larger-augmentation idea.}
Split \(qr\) labels into \(q\) groups of size \(r\). At each node of a complete \(q\)-ary tree, experts in child \(j\)'s subtree predict group \(j\). Some label has learner inclusion probability at most \(k/(qr)=c/q\le1/2\). Reveal such a label and descend into its group. One leaf expert remains perfect, while the learner incurs expected loss at least \(1/2\) per level.

\section{Consequences for learning}
\label{sec:consequences}

\subsection{Finite-class agnostic list learning}

Let \(B_1,\ldots,B_N:X\to\cB_r\) be a finite comparator class and let \(P\) be any distribution on \(X\times\cY\). Write \[ \risk(B_i)=\Pp_{(X,Y)\sim P}\{Y\notin B_i(X)\}. \] The learned predictor is a randomized kernel \(x\mapsto Q_x\in\DeltaS(\cA_k)\), with risk averaged over its internal randomization.

\begin{corollary}[Fast stochastic aggregation]
\label{cor:stochastic}
For every \(k>r\), there is a learner based on \(n\) iid samples such that
\begin{equation}
\label{eq:expected-fast}
\E\risk(\widehat A)
\le\min_{i\in[N]}\risk(B_i)
+\frac{\log N}{\eta_N(r,k)n}.
\end{equation}
There is also a learner which, with probability at least \(1-\delta\), satisfies
\begin{align}
\risk(\widehat A)-\min_i\risk(B_i)
\le {}&
\frac{2e\lceil\log(2/\delta)\rceil\log N}
{\eta_N(r,k)n}
\nonumber\\
&+O\left(
\frac{\log(1/\delta)+\sqrt{\log(1/\delta)\log n}}
{\eta_N(r,k)n}
+\frac{\log(\log n/\delta)}n
\right).
\label{eq:highprob-fast}
\end{align}
\end{corollary}

The expectation bound follows directly by online-to-batch conversion of Theorem~\ref{thm:regret}. The high-probability statement uses the model-selection aggregation theorem of \citet{mehta2017fast}. The surrogate prediction is a convex combination of expert incidence vectors and its loss is \[ \widetilde\ell_\eta(v,y)=b_\eta(v_y). \] It is bounded and \(\eta\)-exp-concave because \[ \exp\{-\eta\widetilde\ell_\eta(v,y)\} =e^{-\eta}+(1-e^{-\eta})v_y \] is affine in \(v\). At every comparator vertex, the surrogate equals omission loss. The resource-augmented substitution then converts the learned convex predictor pointwise to an exact randomized \(k\)-list without increasing risk. Appendix~\ref{app:stochastic} checks the assumptions and constants.

For \(r=1\), Corollary~\ref{cor:stochastic} competes with the best ordinary multiclass classifier while returning a list of size \(k\). It gives expected sample complexity \[ O\left(\frac{\log N}{\eta_N(1,k)\varepsilon}\right), \] rather than the usual \(1/\varepsilon^2\) agnostic scale. This makes precise the hedging mechanism highlighted by \citet{pabbaraju2026optimal}. It is a finite-class result and does not establish a dimension-only \(1/\varepsilon\) theorem for arbitrary infinite classes.

\subsection{List Littlestone covers}

Let \(\cH\subseteq\cB_r^X\) be a class of \(r\)-list predictors on a finite label alphabet of size \(L\). Let \(\mathrm L_r(\cH)=d\) be its \(r\)-list Littlestone dimension in the sense of \citet{moran2023list}. Their list Sauer lemma provides a horizon-\(T\) cover by at most
\begin{equation}
\label{eq:cover-size}
M_T
\le\sum_{i=0}^d{T\choose i}
\left(r\left\lceil\frac{L-r}{r}\right\rceil\right)^i
\end{equation}
adaptive \(r\)-list experts.

\begin{corollary}[Logarithmic horizon remainder]
\label{cor:sequential}
For every \(k>r\) and \(1\le d\le T\), there is an online randomized \(k\)-list learner satisfying
\begin{equation}
\label{eq:sequential}
\E\Reg_T(\cH)
\le\frac{\log M_T}{\eta_{M_T}(r,k)}
\le
\frac d{\eta_\infty(r,k)}
\log\left(
\frac{eT}{d}
r\left\lceil\frac{L-r}{r}\right\rceil
\right).
\end{equation}
\end{corollary}

For any realized sequence and comparator \(h\), change the label only on rounds where \(h\) errs, choosing an arbitrary label from \(h(x_t)\). The modified sequence is realizable by \(h\). A cover expert includes every modified label, so it errs on no more rounds than \(h\) on the original sequence. Aggregating the cover with Theorem~\ref{thm:regret} proves the result.

The prior generic aggregation in \citet{moran2023list} gives a square-root-in-\(T\) remainder. Strictly more learner slots reduce that term to \(\log M_T=O(d\log T)\). For a finite comparator class, using the class itself as the cover removes the horizon dependence entirely.

\section{Related work and limitations}
\label{sec:related}

\paragraph{List prediction.}
\citet{charikar2023characterization} characterize PAC list learnability through list DS dimension. \citet{moran2023list} characterize online list learnability and show that larger learner lists can yield negative regret, but their quantitative resource-augmented upper bound retains a square-root horizon residual. \citet{hanneke2025private} separate private and online list learnability. \citet{kosoy2026ambiguous} studies a stronger ambiguous-prediction condition which requires the predicted set to contain the outcome and exclude every predictably wrong label.

\paragraph{Dual set-valued protocols.}
\citet{pour2026multiple} study a complementary problem where the target is a set of acceptable labels and the learner predicts one label. Their comparator is penalized when its entire predicted set differs from the revealed target set. They obtain constant regret for special multi-label classes under set-valued feedback. In our protocol the target is one label, both sides predict lists, and both receive omission loss. Neither result contains the other.

\paragraph{Mixability and abstention.}
Mixability characterizes constant regret when experts and learner share an action space \citep{vovk1998game,vanerven2012mixability,reid2015generalized}. Resource-augmented mixability instead restricts expert actions and enlarges learner actions. \citet{neu2020abstention} show that a binary reject option with cost below \(1/2\) gives \(\Theta(\log N/(1-2c))\) regret. A fixed-cost abstention action and a fixed-cardinality multiclass set have different geometry, although both exhibit a sharp transition from square-root to constant regret.

\paragraph{Limitations.}
The guarantee controls expected omission loss and requires fresh learner randomization. It does not yield a deterministic exact-list predictor with the same pointwise substitution. The finite-class stochastic result does not by itself resolve agnostic list learning for every class of finite list DS dimension. The large-augmentation lower bound is not tight when \(k/r\) approaches a fixed \(N\), where the exact finite-pool upper bound vanishes. Finally, computational efficiency assumes the expert lists can be enumerated and that filler labels can be accessed.

\bibliography{references}
\bibliographystyle{iclr2027_conference}

\appendix

\section{Proof of the exact geometry}
\label{app:geometry}

\subsection{The hypersimplex criterion}

Fix a finite active alphabet of size \(m\). A distribution \(Q\) over exact \(k\)-sets has inclusion marginals \(p\in[0,1]^m\) with \(\sum_yp_y=k\). Conversely, every such vector is a convex combination of exact \(k\)-set incidence vectors. One proof repeatedly chooses two fractional coordinates and moves mass between them until one becomes integral, randomizing between the two boundary points to preserve the current vector. The process terminates at incidence vectors and expresses the original point as their convex combination.

For a posterior over comparator lists,
\begin{align*}
-\frac1\eta\log\sum_iw_i e^{-\eta\ell(B_i,y)}
&=-\frac1\eta\log\left(s_y+(1-s_y)e^{-\eta}\right)\\
&=b_\eta(s_y).
\end{align*}
The learner's expected loss is \(1-p_y\). Hence the substitution inequalities are \(p_y\ge g_\eta(s_y)\) for all \(y\). A feasible \(p\) can dominate these lower bounds exactly when their sum is at most \(k\), proving Lemma~\ref{lem:budget}.

\subsection{Finite alphabet}

The posterior marginals satisfy \(\sum_ys_y\le r\). Since \(g_\eta\) is increasing and concave,
\begin{equation}
\label{eq:jensen-m}
\sum_{y=1}^m g_\eta(s_y)
\le m g_\eta\left(\frac{\sum_ys_y}{m}\right)
\le m g_\eta(r/m).
\end{equation}
Uniform mixing over all exact \(r\)-subsets gives \(s_y=r/m\), so both inequalities are equalities. By Lemma~\ref{lem:budget}, the exact threshold is \(m g_\eta(r/m)=k\), which is \eqref{eq:finite-m}. If \(k\ge m\), the learner includes the entire alphabet and has zero loss.

For fixed \(s\in(0,1)\), the function \(g_\eta(s)\) is strictly increasing in \(\eta>0\). Indeed, it is the secant slope from zero to \(\eta\) of the strictly convex function \[ \phi_s(\eta)=\log(1-s+se^\eta). \] It increases from \(s\) to one. Since \(r<k<m\), the finite-alphabet threshold has a unique positive root.

\subsection{Fixed expert pool}

A posterior on at most \(N\) experts has no more than \(Nr\) active labels. Applying \eqref{eq:jensen-m} on its active support gives \[ \sum_yg_\eta(s_y) \le Nr\,g_\eta(1/N). \] Equality is attained by uniform mixing over \(N\) pairwise disjoint exact \(r\)-lists. The threshold \(Nr\,g_\eta(1/N)=k\) is equivalent to \eqref{eq:finite-n}. Its left side after division by \(r\) increases continuously from one to \(N\) as \(\eta\) ranges from zero to infinity. If \(k\ge Nr\), the union of all posterior-supported lists fits in the learner action.

\subsection{Alphabet-uniform constant and monotonicity}

Concavity and \(g_\eta(0)=0\) give the tangent bound \[ g_\eta(s)\le g_\eta'(0)s =\frac{e^\eta-1}{\eta}s. \] Therefore \[ \sum_yg_\eta(s_y) \le r\frac{e^\eta-1}{\eta}. \] This is at most \(k\) exactly under \eqref{eq:infinite}. Conversely, the uniform posterior over all \(r\)-subsets of an \(m\)-label alphabet satisfies \[ m g_\eta(r/m)\longrightarrow r g_\eta'(0) =r\frac{e^\eta-1}{\eta}. \] Any larger \(\eta\) fails on all sufficiently large alphabets.

For a concave function vanishing at zero, \(g_\eta(s)/s\) decreases in \(s\). Hence \[ m g_\eta(r/m) =r\frac{g_\eta(r/m)}{r/m} \] increases with \(m\). The corresponding root decreases to \(\eta_\infty\). The same argument applied to \(N g_\eta(1/N)\) proves the fixed-pool monotonicity.

\section{Exact-cardinality substitution}
\label{app:rounding}

We give a direct rounding procedure which avoids enumerating the vertices of the hypersimplex. Start with demands \(d_y=g_\eta(s_y)\). By Lemma~\ref{lem:budget}, their sum is at most \(k\). Increase coordinates in any order without exceeding one until obtaining \(p\ge d\) with \(\sum_yp_y=k\). There is sufficient capacity because \(|\cY|\ge k\).

\begin{minipage}{0.94\linewidth}
\textbf{Algorithm 1: systematic exact-list rounding}
\begin{enumerate}
\item Fix an ordering \(y_1,\ldots,y_m\) containing every positive coordinate and enough filler labels. Put \(C_0=0\) and \(C_j=\sum_{i=1}^jp_{y_i}\).
\item Draw \(U\) uniformly from \([0,1)\).
\item Return every \(y_j\) for which the half-open interval \([C_{j-1},C_j)\) contains a point in
\[ \{U,U+1,\ldots,U+k-1\}. \]
\end{enumerate}
\end{minipage}

\begin{lemma}[Rounding guarantee]
\label{lem:rounding}
Algorithm~1 returns exactly \(k\) distinct labels and satisfies \[ \Pp(y\text{ is returned})=p_y \quad\text{for every }y. \]
\end{lemma}

\begin{proof}
The intervals partition \([0,k)\), which contains exactly the \(k\) displayed points. Each interval has length \(p_y\le1\), so a half-open interval cannot contain two points separated by one. The returned labels are therefore distinct and number exactly \(k\).

For a fixed interval \(I=[C_{j-1},C_j)\), label \(y_j\) is selected when \(U\) belongs to the projection of \(I\) modulo one. Because \(I\) has length at most one, the measure of this projection counted over the integer translates is exactly \(|I|=p_{y_j}\). Uniformity of \(U\) gives the marginal claim.
\end{proof}

The construction takes linear time in the ordered support. It can be streamed by advancing through cumulative sums and the \(k\) grid points simultaneously.

\section{Online regret proof}
\label{app:online}

Let \(W_t=\sum_iW_{t,i}\). The exponential update gives
\begin{align}
\frac{W_{t+1}}{W_t}
&=\sum_iw_{t,i}e^{-\eta\ell(B_{t,i},y_t)}\nonumber\\
&=e^{-\eta}+(1-e^{-\eta})s_{t,y_t}.
\label{eq:potential-step}
\end{align}
By the substitution and \eqref{eq:bg}, \[ \E_t\ell(A_t,y_t) \le b_\eta(s_{t,y_t}) =-\frac1\eta\log\frac{W_{t+1}}{W_t}. \] Summing over rounds yields \[ \sum_{t=1}^T\E_t\ell(A_t,y_t) \le-\frac1\eta\log\frac{W_{T+1}}{W_1}. \] Initialize \(W_{1,i}=\pi_i\), so \(W_1=1\). For every \(i\), \[ W_{T+1} \ge\pi_i\exp\left\{-\eta\sum_{t=1}^T \ell(B_{t,i},y_t)\right\}. \] Substitution proves Theorem~\ref{thm:regret}.

\section{Lower-bound proofs}
\label{app:lower-proofs}

\subsection{A discounted binary embedding}

Let \(k=r+s\) and \(a=s/r<1\). On the alphabet \([r]\times\{0,1\}\), embed a binary prediction \(h(x)\) as \[ B_h(x)=\{(j,h(x)):j\in[r]\}. \] For any distribution over exact \(k\)-lists, let \(p_{j,b}\) denote the inclusion marginal of \((j,b)\) and define \[ \bar p_b=\frac1r\sum_{j=1}^rp_{j,b}. \] Since the total inclusion budget is \(r+s\), \[ \bar p_0+\bar p_1=1+a. \] Each \(\bar p_b\le1\), so \(\bar p_b\ge a\). Therefore \[ u_b=\frac{\bar p_b-a}{1-a} \] is a probability distribution on \(\{0,1\}\).

Draw \(J\) uniformly from \([r]\). For either binary outcome \(Y\), \[ \E_J\ell(A,(J,Y)) =1-\bar p_Y =(1-a)(1-u_Y). \] Every embedded expert has \[ \ell(B_h(x),(J,Y))=\1\{h(x)\ne Y\}. \] This proves the exact discount identity \eqref{eq:discount}. When \(a=0\), it transfers every ordinary binary experts lower bound, proving the first part of Theorem~\ref{thm:lower}.

\subsection{Noisy hypercube near the threshold}

Let \(d=\lfloor\log_2N\rfloor\). It suffices to lower bound a pool of \(2^d\) experts because duplicates can be added. Index experts by \(\theta\in\{0,1\}^d\) and let \(h_\theta(x)=\theta_x\) on \(X=[d]\).

We use Yao's principle. Draw \(\theta\) uniformly. Independently on every round, draw \(X_t\) uniformly from \([d]\) and \(J_t\) uniformly from \([r]\). Conditional on \(X_t\) and \(\theta\), draw \[ \Pp(Y_t=\theta_{X_t})=\frac{1+\gamma}{2}, \qquad \gamma=4a. \] The revealed multiclass outcome is \((J_t,Y_t)\). We lower bound every deterministic list strategy against this distribution. This also covers randomized strategies after averaging their internal seed.

Fix the current coordinate \(x=X_t\). Let \(P_0\) and \(P_1\) be the laws of the history through round \(t-1\), conditional on \(\theta_x=0\) and \(\theta_x=1\), with the other latent bits averaged uniformly. A previous observation carries information about \(\theta_x\) only when its coordinate equals \(x\). For \(\gamma\le1/4\),
\begin{align}
\KL\left(
\Ber\left(\frac{1+\gamma}{2}\right)
\middle\|
\Ber\left(\frac{1-\gamma}{2}\right)
\right)
&=\gamma\log\frac{1+\gamma}{1-\gamma}\nonumber\\
&\le4\gamma^2.
\label{eq:ber-kl}
\end{align}
The chain rule and the \(1/d\) probability of querying \(x\) give \[ \KL(P_0\|P_1) \le\frac{4\gamma^2(t-1)}d. \] If \(t-1\le d/(32\gamma^2)\), then Pinsker's inequality gives \[ \TV(P_0,P_1)\le\sqrt{\KL(P_0\|P_1)/2}\le\frac14. \] The Bayes error for predicting the uniform bit \(\theta_x\) from the history is at least \[ \frac{1-\TV(P_0,P_1)}2\ge\frac38. \] The same lower bound applies to the induced randomized binary prediction \(u\). Let \(q_t=\E[1-u_{\theta_{X_t}}]\). Its expected binary error on the noisy label is \[ \frac{1-\gamma}{2}+\gamma q_t. \] The true expert \(h_\theta\) has expected error \((1-\gamma)/2\). Using the discount identity and \(q_t\ge3/8\), the expected regret relative to \(h_\theta\) on round \(t\) is at least
\begin{align*}
(1-a)\left(\frac{1-\gamma}{2}+\gamma q_t\right)
-\frac{1-\gamma}{2}
&\ge(1-a)\gamma\frac38-a\frac{1-\gamma}{2}\\
&=a+\frac{a^2}{2}\\
&\ge a.
\end{align*}
Set \[ T=1+\left\lfloor\frac{d}{32\gamma^2}\right\rfloor. \] Every round through \(T\) satisfies the information bound, and \(T>d/(32\gamma^2)\). Expected regret relative to \(h_\theta\), hence also relative to the best expert, is at least \[ aT>\frac{ad}{32\gamma^2} =\frac{d}{512a}. \] This proves \eqref{eq:near-lower} with an explicit universal constant.

\subsection{Grouped tree for larger augmentation}

Let \(c=k/r\), \(q=\lceil2c\rceil\), and \(d=\lfloor\log_qN\rfloor\). Use \(qr\) labels partitioned into disjoint groups \(G_1,\ldots,G_q\), each of size \(r\). Consider a complete \(q\)-ary tree of depth \(d\), with one expert for each leaf. At an internal node, every expert in child \(j\)'s subtree predicts \(G_j\). Predictions of experts already outside the retained subtree can be arbitrary.

At the current node, let \(p_y\) be the learner's inclusion marginals. Since \(\sum_yp_y=k\), some label \(y\) satisfies \[ p_y\le\frac{k}{qr}=\frac cq\le\frac12. \] The adversary reveals \(y\) and descends to the child indexed by the group containing \(y\). The learner's expected loss is at least \(1/2\). Every expert retained after this step predicted the whole group containing \(y\), so it has zero loss. Repeating for \(d\) levels leaves one zero-loss expert and gives regret at least \(d/2\).

If \(N\) is not a power of \(q\), add arbitrary duplicate experts to the \(q^d\)-expert construction. When \(N\ge q^2\), elementary comparisons between \(q=\lceil2c\rceil\) and \(c+1\) give \[ \frac d2 =\Omega\left(\frac{\log N}{\log(c+1)}\right). \] This proves \eqref{eq:large-lower}.

\section{Stochastic conversions}
\label{app:stochastic}

\subsection{Expected excess risk}

Run the online algorithm on an iid sample \(Z_1,\ldots,Z_n\), with \(Z_t=(X_t,Y_t)\). Let \(Q_t(\cdot\mid x)\) be the randomized list predictor formed from the weights before observing \(Z_t\). After training, draw \(J\) uniformly from \([n]\) and output \(Q_J\). The prediction \(Q_t\) depends only on \(Z_{<t}\), so independence gives \[ \E\ell(Q_t(X_t),Y_t)=\E\risk(Q_t). \] Theorem~\ref{thm:regret} and the inequality \(\E\min_i V_i\le\min_i\E V_i\) imply
\begin{align*}
\sum_{t=1}^n\E\risk(Q_t)
&\le
\E\min_i\sum_{t=1}^n\ell(B_i(X_t),Y_t)
+\frac{\log N}{\eta_N(r,k)}\\
&\le
n\min_i\risk(B_i)
+\frac{\log N}{\eta_N(r,k)}.
\end{align*}
Dividing by \(n\) proves \eqref{eq:expected-fast}.

\subsection{High-probability excess risk}

Map every expert output to its incidence vector \(v_i(x)\in\{0,1\}^{\cY}\). The action space for the surrogate is the convex hull of the \(N\) vector-valued functions. Define \[ \widetilde\ell_\eta(v,(x,y))=b_\eta(v(x)_y). \] The loss lies in \([0,1]\), and its negative exponential is affine: \[ e^{-\eta\widetilde\ell_\eta(v,(x,y))} =e^{-\eta}+(1-e^{-\eta})v(x)_y. \] Thus it is \(\eta\)-exp-concave. At every expert vertex, \[ \widetilde\ell_\eta(v_i,(x,y)) =\1\{y\notin B_i(x)\}. \]

Apply Theorem 4 of \citet{mehta2017fast} with the uniform prior and diameter bound \(B=1\). Its PM-EWOO estimator \(\widehat v\) satisfies, with probability at least \(1-\delta\),
\begin{align*}
\widetilde\risk(\widehat v)-\min_i\risk(B_i)
\le {}&
\frac{2e\lceil\log(2/\delta)\rceil\log N}{\eta n}\\
&+O\left(
\frac{\log(1/\delta)+\sqrt{\log(1/\delta)\log n}}{\eta n}
+\frac{\log(\log n/\delta)}n
\right).
\end{align*}
For every \(x\), \(\widehat v(x)\) is a mixture of the \(N\) expert incidence vectors. The fixed-pool budget theorem therefore supplies a randomized exact \(k\)-list whose conditional omission probability at \(y\) is at most \(b_\eta(\widehat v(x)_y)\). Taking \(\eta=\eta_N(r,k)\) proves \eqref{eq:highprob-fast}.

\section{Sequential-cover conversion}
\label{app:sequential}

Let \(\mathcal F_T\) be the \(T\)-cover supplied by the list Sauer lemma of \citet{moran2023list}, with size bounded by \eqref{eq:cover-size}. Fix an arbitrary realized sequence \((x_t,y_t)_{t=1}^T\) and a comparator \(h\in\cH\). Construct a modified label sequence by \[ y'_t=
\begin{cases}
y_t,&y_t\in h(x_t),\\
\text{an arbitrary element of }h(x_t),&y_t\notin h(x_t).
\end{cases}
\] The modified sequence is realizable by \(h\). Hence some adaptive \(r\)-list expert \(f\in\mathcal F_T\) includes \(y'_t\) on every round. Whenever \(h\) is correct on the original sequence, \(y'_t=y_t\), so \(f\) is also correct. Therefore \[ \sum_{t=1}^T\ell(f_t,y_t) \le\sum_{t=1}^T\ell(h(x_t),y_t). \] Aggregate the \(M_T\) cover experts with Theorem~\ref{thm:regret}. This gives the first inequality in \eqref{eq:sequential}. For the second, use \(\eta_{M_T}\ge\eta_\infty\) and \[ \sum_{i=0}^d{T\choose i}a^i \le a^d\sum_{i=0}^d{T\choose i} \le\left(\frac{eTa}{d}\right)^d \] for \(a=r\lceil(L-r)/r\rceil\).

\section{Asymptotic calculations}
\label{app:asymptotics}

Write \(c=k/r\). The equation for \(\eta_\infty\) is \[ \frac{e^\eta-1}{\eta}=c. \] As \(\eta\downarrow0\), \[ \frac{e^\eta-1}{\eta} =1+\frac\eta2+\frac{\eta^2}{6}+O(\eta^3), \] which gives \(\eta=2(c-1)+O((c-1)^2)\). As \(c\to\infty\), the equation is \(e^\eta=1+c\eta\). Taking logarithms and one fixed-point iteration gives \[ \eta=\log c+\log\log c+o(1). \]

For fixed \(N\), let \(c\uparrow N\). The finite-pool equation obeys \[ \frac N\eta\log\left(1+\frac{e^\eta-1}{N}\right) =N-\frac{N\log N}{\eta}+o(1/\eta). \] Therefore \[ \eta_N(r,k)\sim\frac{N\log N}{N-c}. \]

\section{Numerical verification scope}
\label{app:verification}

We independently checked the proof identities in finite instances. The checks enumerate all exact \(k\)-sets for small alphabets and solve the substitution feasibility problem as a linear program. They verify feasibility immediately below each claimed finite-alphabet threshold and infeasibility immediately above it. Random posteriors over all \(r\)-sets satisfy the budget inequality, while the uniform posterior attains equality. Additional tests verify monotonic convergence of the finite-\(m\) and finite-\(N\) constants, the exact discount identity, the systematic-rounding marginals, and the online potential bound over simulated adaptive sequences. These checks support the algebra but are not used as proof.

\end{document}
